\documentclass{article}

\usepackage[main]{neurips_2026}
%  \usepackage[main]{neurips_2026}
% \usepackage[preprint]{neurips_2026}

\usepackage[utf8]{inputenc}
\usepackage[T1]{fontenc}
\usepackage{hyperref}
\usepackage{url}
\usepackage{booktabs}
\usepackage{amsfonts}
\usepackage{amsmath}
\usepackage{nicefrac}
\usepackage{microtype}
\usepackage{xcolor}
\usepackage{graphicx}
\usepackage{algorithm}
\usepackage{algpseudocode}
\usepackage{multirow}

% =============================================================================
% TITLE 候选 (挑一个或自己改写)
%   1. Function-Aware Fill-in-the-Middle as Mid-Training for Coding Agents
%   2. Bridging Pretraining and Agentic Post-Training:
%      Structured FIM Mid-Training for Coding Agent Foundation Models
%   3. Structure-Aware FIM: A Scalable Mid-Training Recipe for
%      Coding Agents from Self-Supervised Code
% 写作建议: 标题里同时点出 (a) FIM, (b) mid-training, (c) coding agent
% =============================================================================

\title{Function-Aware Fill-in-the-Middle as Mid-Training\\for Coding Agent Foundation Models}

\author{%
  Anonymous Authors \\
  Anonymous Institution \\
  \texttt{anonymous@example.com}
}


\begin{document}
\maketitle

% =============================================================================
% ABSTRACT  (~ 200 words, 1 段)
% =============================================================================
% 写作模板 (按这个顺序写, 每点 1-2 句):
%   (1) 背景: coding agents 需要在 tool return 之后继续推理; 这种结构
%       在 standard next-token pretraining 中暴露不充分.
%   (2) 关键观察: function call dependencies 在结构上与 agent 的
%       (action -> external return -> continuation) 同构, 且在 GitHub
%       上以巨大规模天然存在.
%   (3) 方法: function-aware FIM mid-training. PDG + 信息论打分选择
%       mask 的函数; Gemini-3 生成 CoT.
%   (4) 实验: 在三轴 robustness 上验证:
%       (a) 跨 model size (Qwen2.5-Coder 7B/14B/32B)
%       (b) 跨 post-training pipeline (r2e-gym, swe-smith)
%       (c) 跨 base model family (Qwen2.5-Coder, Qwen3-8B)
%   (5) 结果: SWE-Bench-Verified 上 7B +X.X / 14B +Y.Y / 32B +Z.Z;
%       compute-matched 下显著优于扩充 post-training 数据.
%   (6) Cross-domain transfer: 仅用 Python 代码训练, 在 non-coding
%       tool-use benchmark (tau-bench, BFCL) 上也观察到提升, 直接验证
%       "function call ↔ tool call 同构" 的 hypothesis.
%   (7) Takeaway: structured FIM mid-training 是一种 scalable,
%       self-supervised 的 agent capability prior.
% CAUTION: cross-base-model claim 措辞要谨慎, 用 "we further validate on a
%          different base family (Qwen3-8B), suggesting the gain generalizes
%          beyond Qwen2.5-Coder", 不要说 "works on any base model".
% =============================================================================
\begin{abstract}
% TODO: 完整撰写 abstract, 严格 200 词以内, 一段
Coding agents must reason over tool returns and continue execution under
external perturbations---a capability poorly exposed by standard next-token
pretraining on code. We observe that \emph{function call dependencies} in
real-world code are structurally isomorphic to the
\textit{action $\rightarrow$ tool-return $\rightarrow$ continuation} loop of
coding agents, and exist at internet scale on GitHub. Building on this, we
propose \textbf{function-aware fill-in-the-middle (FIM) mid-training}: a
self-supervised objective that masks selected functions whose call sites are
identified via program dependency graph analysis and an information-theoretic
predictability score, augmented with chain-of-thought rationales generated by
a frontier model.
% [continue with experimental setup across three axes,
%  headline numbers,
%  cross-domain transfer to tau-bench/BFCL highlight,
%  takeaway]
\end{abstract}


% =============================================================================
% 1. INTRODUCTION  (~ 1.25 页)
% =============================================================================
% 结构 (5 段):
%   段 1: Coding agent 已成为重要应用; 但能力主要靠堆 post-training 数据.
%         pretraining --> [GAP] --> post-training. 我们关心这个 gap.
%   段 2: 关键观察 (paper 的 hook).
%         agent loop 与 function call site 结构同构. standard code
%         pretraining 的单向 next-token loss 只暴露了"forward direction".
%         FIM 提供双向 conditioning.
%   段 3: random-span FIM 已经被 base model (Qwen-Coder) 用过, 直接应用不够.
%         我们的 design choice: function-aware (PDG + 信息论打分) + CoT 监督
%         + mid-training stage 定位.
%   段 4: 实验概览: 三轴 robustness (跨 size, 跨 pipeline, 跨 base family);
%         核心数字: 14B 在 SWE-Bench-Verified 上 26.20% -> 29.20%;
%         compute-matched 下显著优于扩充 post-training 数据;
%         cross-domain transfer 到 tau-bench / BFCL.
%   段 5: Contributions (5 条):
%         - Conceptual framing: FIM ↔ agent inductive bias 的同构论证.
%         - Method: function-aware FIM with PDG + IT scoring + CoT.
%         - Empirical robustness across three axes:
%           (a) sizes 7B/14B/32B,
%           (b) pipelines r2e-gym/swe-smith,
%           (c) base families Qwen2.5-Coder/Qwen3.
%         - Direct motivation validation: tau-bench/BFCL transfer.
%         - 数据: 990 decontaminated repo 的 mid-training corpus.
%
% FIG 1 [必须]: 第一页 teaser
%   左:  code snippet 高亮 callsite (caller / call / use of return)
%   中:  agent trajectory 高亮 (think / tool_call / observation / continue)
%   右:  bar chart 关键 result (Qwen2.5-Coder-14B baseline vs +ours)
%
% CAUTION:
%   - 不要 over-claim "FIM = agent loop". 用 "structurally isomorphic" /
%     "shares conditioning structure".
%   - Cross-base-model 措辞: "further validate on a different base family",
%     不要说 "works universally".
%   - Contributions 必须有 method-level 的, 不能只是 "we show FIM works".
% =============================================================================
\section{Introduction}
\label{sec:intro}

% TODO 段 1: motivate coding agent + pretraining/post-training gap.
% TODO 段 2: 核心 motivation (FIM ↔ agent loop 结构同构).
% TODO 段 3: 我们 design choice 与 vanilla FIM 的区别.
% TODO 段 4: 实验 setup 一句 (强调三轴 robustness) + headline number.
% TODO 段 5: Contributions (5 条, "三轴 robustness"作为亮点).

% TODO FIG 1: teaser figure


% =============================================================================
% 2. RELATED WORK  (~ 0.6 页)
% =============================================================================
% 四块:
%   2.1 Mid-training / continued pretraining
%       - 引用 MiniCPM, OLMo, Code-Llama long-context mid-training 等
%       - 我们的差异化: agent-oriented mid-training, 用 structured FIM.
%
%   2.2 Fill-in-the-middle pretraining
%       - 引用 Bavarian 2022, Code-Llama, StarCoder, Qwen-Coder, DeepSeek-Coder
%       - 关键: 这些都是 random span FIM, 在 pretraining 阶段大规模混合.
%         我们的差异化: function-aware selection + CoT + mid-training stage.
%   CAUTION: 不要否认前人贡献, frame 成 "extension/specialization".
%
%   2.3 Coding agent foundation models
%       - SWE-Bench, r2e-gym, swe-smith, SWE-agent, OpenHands
%       - Synthetic agent trajectory works (SWE-Gym, Agent-FLAN)
%       - 我们的差异化: 不依赖 trajectory, 用 self-supervised code 信号.
%
%   2.4 Distillation from frontier models [必须有, 主动 disarm]
%       - acknowledge 我们用 Gemini-3 生成 CoT, 引用 Distill-Whisper, Orca,
%         Self-Instruct 等
%       - 强调我们 ablation 拆开 distillation 与 FIM 结构两部分贡献:
%         "We separate the contribution of FIM structure from CoT distillation
%          via a controlled ablation in Section X.X."
% CAUTION: 这一节存在 = 主动 disarm reviewer 的 distillation 质疑.
% =============================================================================
\section{Related Work}
\label{sec:related}

\paragraph{Mid-training and continued pretraining.}
% TODO

\paragraph{Fill-in-the-middle objectives.}
% TODO

\paragraph{Coding agent foundation models.}
% TODO

\paragraph{Distillation from frontier models.}
% TODO (主动 disarm reviewer)


% =============================================================================
% 3. METHOD  (~ 2 页)
% =============================================================================
\section{Method}
\label{sec:method}


% -----------------------------------------------------------------------------
% 3.1 Motivation: Function Call Dependency as Agent Inductive Bias
% -----------------------------------------------------------------------------
% (~ 0.4 页, 1-2 段 + 1 张图)
%
% 段 1: 形式化描述同构关系.
%   agent loop: pi(a_t | h_t), env(o_{t+1} | h_t, a_t),
%               pi(continuation | h_t, a_t, o_{t+1}).
%   function call site: caller_context -> f(args) -> return_value
%                       -> downstream_use(return_value).
%   关键 disclaimer: FIM 是 P(M | prefix, suffix) (suffix 已知);
%        agent inference 时 continuation 是 generated. 我们 hypothesize
%        FIM 训练 induces 一种 representation, 让模型在 inference 时也能
%        利用这种 dependency structure.
%   CAUTION: 这一段是 paper 的 conceptual heart. 一定要写得精确.
%
% FIG 2 [必须]: side-by-side 同构示意图
%   (a) function call site 4 部分: pre-call | call | return | post-use
%   (b) agent step 4 部分:        history   | action | obs    | continuation
% -----------------------------------------------------------------------------
\subsection{Motivation: Function Calls as Agent-Like Structures}
\label{sec:method:motivation}

% TODO: 形式化 isomorphism
% TODO FIG 2: 同构示意图


% -----------------------------------------------------------------------------
% 3.2 Data Collection and Decontamination
% -----------------------------------------------------------------------------
% (~ 0.3 页)
%
% 内容:
%   - 起点: GitHub repo, 通过 star + 10 个 category keyword 收集 2000 个
%   - 人工筛选高质量 1000 个
%   - 去除 SWE-Bench source repo 及相关 fork --> 990 repo
%   - 报告: 总文件数, 总 Python LoC, license 分布
%   - 时间戳: commit cutoff < SWE-Bench task base commit
%
% FIG 3 [必须]: 10 类 category 的 repo 数量分布
% TAB 1 [建议]: 数据集 statistics
%
% CAUTION:
%   - "为什么 990 而不是 999?" --> 说清楚 SWE-Bench 去重后剩多少
%   - "category 是 emergent 还是 designed?" --> 写清楚
%   - 必须明确 repo list 会 release
% -----------------------------------------------------------------------------
\subsection{Data Collection and Decontamination}
\label{sec:method:data}

% TODO: collection + filtering + decontamination
% TODO FIG 3: category distribution
% TODO TAB 1: dataset statistics


% -----------------------------------------------------------------------------
% 3.3 Function Selection via PDG and Information-Theoretic Scoring
% -----------------------------------------------------------------------------
% (~ 0.6 页)
%
% Step 1: PDG 构建 (AST -> function nodes -> caller-callee edges)
% Step 2: Mask 候选 (单 function 或 connected k-tuple)
% Step 3: 信息论打分
%   S(f | context) = -log P_M_ref(f | context_outside_f)
%   保留 S_lo <= S(f) <= S_hi 的 function (sweet spot).
% Step 4: Mask & format (用 Qwen-Coder 原生 FIM tokens)
%
% ALG 1 [建议]: selection pipeline pseudocode
% FIG 4 [可选]: PDG 例子
% -----------------------------------------------------------------------------
\subsection{Function Selection via Dependency Graph and Predictability}
\label{sec:method:selection}

% TODO: 写 step 1-4
% TODO ALG 1: pseudocode


% -----------------------------------------------------------------------------
% 3.4 Chain-of-Thought Augmentation
% -----------------------------------------------------------------------------
% (~ 0.4 页)
%
% 关键 design rationale (回应 distillation 质疑):
%   CoT 不是为了直接 distill, 而是 (a) 显式建模 "从 context 推断 callee
%   行为" 的中间过程, (b) 在 mid-training 阶段保留 reasoning trace 与
%   最终代码的 alignment.
%   --> forward-reference 后面的 ablation (Section 4.7).
%
% FIG 5 [建议]: 一个具体的 training example
% -----------------------------------------------------------------------------
\subsection{Chain-of-Thought Augmentation}
\label{sec:method:cot}

% TODO: CoT 嵌入方式 + design rationale
% TODO FIG 5: 训练样本示例


% -----------------------------------------------------------------------------
% 3.5 Mid-Training and Post-Training Setup
% -----------------------------------------------------------------------------
% (~ 0.3 页)
% TAB 2 [建议, 可放 appendix]: 超参表格
% -----------------------------------------------------------------------------
\subsection{Mid-Training and Post-Training Setup}
\label{sec:method:setup}

% TODO: base model + token budget + hyperparams + post-training pipeline


% =============================================================================
% 4. EXPERIMENTS  (~ 3.4 页)
% =============================================================================
\section{Experiments}
\label{sec:exp}


% -----------------------------------------------------------------------------
% 4.1 Experimental Setup  (~ 0.3 页)
% -----------------------------------------------------------------------------
% 内容:
%   - Benchmarks 拆两类:
%     * Agent benchmarks: SWE-Bench-Verified, SWE-Bench-Lite, tau-bench, BFCL
%     * General coding benchmarks: LiveCodeBench, OJBench, FullStackBench-EN,
%       FullStackBench-ZH, Terminal-Bench, HumanEval+
%   - Eval scaffold: agent harness, max turns, retrieval setting
%   - Statistical protocol: 3 seeds, mean ± std, paired bootstrap
%   - Baselines:
%     * Qwen2.5-Coder-Instruct (no mid-train, no post-train)
%     * r2e-gym (ours-reproduced + official)
%     * swe-smith (ours-reproduced + official)
%     * +FIM-Midtrain + r2e-gym
%     * +FIM-Midtrain + swe-smith
% -----------------------------------------------------------------------------
\subsection{Setup}
\label{sec:exp:setup}

% TODO: benchmarks (两类) + eval scaffold + statistical protocol + baselines


% -----------------------------------------------------------------------------
% 4.2 Main Results: Agent Benchmarks  (~ 0.7 页 + 大表)
% -----------------------------------------------------------------------------
% TAB 3 [必须, paper 主表 1]: Agent benchmark 主结果
%   行: 跨 7B/14B/32B 用 \multirow 分组, 每组内:
%     - Base instruct
%     - r2e-gym (ours-reproduced)
%     - r2e-gym (official report)         [灰色, "reported"]
%     - +FIM-Midtrain + r2e-gym           [bold, ours]
%     - swe-smith (ours-reproduced)
%     - swe-smith (official report)       [灰色]
%     - +FIM-Midtrain + swe-smith         [bold, ours]
%   列: SWE-Bench-V, SWE-Bench-L, tau-bench, BFCL, average
%   每个数字带 ± std (3 seeds).
%
% 写作 critical paragraph (paper 最强卖点之一):
%   单独一段 highlight tau-bench / BFCL 的提升:
%   "Crucially, our mid-training corpus contains exclusively Python code
%    with no tool-use trajectory data. Yet we observe non-trivial
%    improvements on tau-bench (+X.X) and BFCL (+Y.Y), suggesting that the
%    inductive bias learned from function call structures transfers to
%    general tool-use scenarios. This directly supports our hypothesis
%    that function calls and agent tool calls share an underlying
%    conditioning structure."
% -----------------------------------------------------------------------------
\subsection{Main Results on Agent Benchmarks}
\label{sec:exp:main}

% TODO TAB 3: 主结果表 1 (agent: SWE-V/L, tau-bench, BFCL)
% TODO: 主结果段落
% TODO: 单独一段 highlight cross-domain transfer to tau-bench/BFCL
%       (motivation 直接验证, paper 最强 selling point)


% -----------------------------------------------------------------------------
% 4.3 Generalization to General Coding Benchmarks  (~ 0.5 页 + 表)
% -----------------------------------------------------------------------------
% TAB 4 [必须, paper 主表 2]: General coding benchmark 表
%   行结构同 TAB 3
%   列: LiveCodeBench, OJBench, FullStackBench-EN, FullStackBench-ZH,
%       Terminal-Bench, HumanEval+, average
%
% 写作 takeaway:
%   - 不仅没掉分, 还有持续提升
%   - Terminal-Bench 提升 (closest 到 non-coding tool-use)
%   - FullStackBench-ZH 上的 transfer (虽训练数据是 Python, 多语言也有效)
% -----------------------------------------------------------------------------
\subsection{Generalization to General Coding Benchmarks}
\label{sec:exp:gen}

% TODO TAB 4: 主结果表 2 (general coding)
% TODO: generalization narrative


% -----------------------------------------------------------------------------
% 4.4 Cross-Base-Model Validation: Qwen3-8B  (~ 0.4 页 + 小表)  [新增]
% -----------------------------------------------------------------------------
% 这一节支持 abstract/intro 的 cross-base-model claim, 不能埋在 ablation.
%
% TAB 5 [必须]: Qwen3-8B 小表
%   行: Qwen3-8B-Base / +r2e-gym / +FIM-Midtrain + r2e-gym
%   列: SWE-Bench-V, SWE-Bench-L, average
%
% 写作要点:
%   - "to verify the approach is not specific to Qwen2.5-Coder, we repeat
%     the core experiment on Qwen3-8B"
%   - 强调一致性: gain magnitude vs Qwen2.5-Coder-7B 是否接近?
%   - 诚实承认: due to compute constraints, we evaluate on SWE-Bench
%     subsets only; full-benchmark validation on Qwen3-8B is left to
%     future work.
%
% CAUTION:
%   - 措辞谨慎, 一组 base model 不能 generalize 到 "all base model"
%   - 设置必须和 Qwen2.5-Coder 一致 (token budget, 超参), 否则不 fair
% -----------------------------------------------------------------------------
\subsection{Cross-Base-Model Validation}
\label{sec:exp:crossbase}

% TODO TAB 5: Qwen3-8B 小表
% TODO: cross-base-model 段落 (措辞谨慎)
% TODO 实验: Qwen3-8B + r2e-gym 的 baseline 与 +FIM 对比 (在跑)


% -----------------------------------------------------------------------------
% 4.5 Scaling Across Model Sizes  (~ 0.3 页 + 1 张图)
% -----------------------------------------------------------------------------
% FIG 6 [必须]: scaling curve. X 轴 model size, Y 轴 SWE-Bench-Verified.
%   两条线: r2e-gym only, +FIM-Midtrain.
%
% 写作 framing (备 plan A 和 plan B):
%   Plan A (gain 稳定/扩大): "method scales".
%   Plan B (gain 缩小): "更大 model 已从 pretraining 学到部分结构 prior,
%       边际收益递减, 但在 7B/14B 这个 practical deployment range
%       价值显著."
% -----------------------------------------------------------------------------
\subsection{Scaling Across Model Sizes}
\label{sec:exp:scaling}

% TODO FIG 6: scaling curve
% TODO: 等 32B 跑完后写 narrative


% -----------------------------------------------------------------------------
% 4.6 Compute-Matched Comparison  (~ 0.4 页 + 1 张图)
% -----------------------------------------------------------------------------
% FIG 7 [必须]: 双线曲线
%   X 轴 = total cost ($1k, $2k, $3k API credit)
%   Y 轴 = SWE-Bench-Verified
%   线 A: 全部花在 post-training data 扩充
%   线 B: 全部花在 mid-training (or mix)
% -----------------------------------------------------------------------------
\subsection{Mid-Training vs Post-Training Under Equal Budget}
\label{sec:exp:budget}

% TODO FIG 7: cost-vs-perf curve


% -----------------------------------------------------------------------------
% 4.7 Ablation Studies  (~ 0.7 页, 表 + 文字 + 图)
% -----------------------------------------------------------------------------
% Ablation 维度 (全部在 14B + r2e-gym 上 controlled):
%
%   (A) CoT source [critical, disarm distillation 质疑]:
%       - Baseline (no mid-train)
%       - FIM mid-train, no CoT
%       - FIM mid-train, 14B-self CoT       (Qwen2.5-Coder-14B-Instruct
%                                            自己生成的 CoT)
%       - FIM mid-train, Gemini-3 CoT       (full method)
%       理想 narrative:
%         "Removing CoT reduces SWE-Bench-V from X to Y but still beats
%          no mid-training by Z, indicating FIM structure contributes
%          independently of distillation signal. Replacing Gemini-3 with
%          self-generated CoT recovers most but not all of the gain,
%          suggesting that CoT quality matters but is not the sole driver."
%       CAUTION: 主动写 "due to compute constraints, we conduct this
%                ablation on the 14B configuration".
%
%   (B) Selection algorithm:
%       - Random function selection
%       - PDG only (no IT scoring)
%       - IT scoring only (random function + IT filter)
%       - Full (PDG + IT)
%
%   (C) Mask granularity (1 / 2 / 3 connected functions / mixed)
%
%   (D) IT score threshold sensitivity (S_lo, S_hi sweep)
%
%   (E) [新增, 强烈建议] Mid-training token budget scaling
%       - 1B / 5B / 10B / 20B mid-training tokens 的 SWE-Bench-V 曲线
%       - 回答 "more is better?", 揭示 saturation / forgetting
%       FIG 9 [建议]: token budget scaling curve
%
% TAB 6 [必须]: 综合 ablation 表 (controlled on 14B + r2e-gym)
%
% 写作: 每个 ablation 给一句 takeaway.
% -----------------------------------------------------------------------------
\subsection{Ablation Studies}
\label{sec:exp:ablation}

% TODO TAB 6: 综合 ablation 表
% TODO 实验 (待补): CoT source ablation on 14B (4 配置)
% TODO 实验 (强烈建议): mid-training token budget scaling
% TODO FIG 9: token budget scaling curve


% =============================================================================
% 5. ANALYSIS  (~ 0.9 页)
% =============================================================================
\section{Analysis}
\label{sec:analysis}


% -----------------------------------------------------------------------------
% 5.1 Behavioral Analysis  (~ 0.4 页)
% -----------------------------------------------------------------------------
% Trajectory-level metrics, mid-training 前后对比:
%   - Average tool calls per task
%   - Recovery rate after error / unexpected output
%   - Average context usage at success
%   - max-turn 限制下的 success curve
%
% TAB 7 [建议]: trajectory metrics 对比
% -----------------------------------------------------------------------------
\subsection{Behavioral Changes in Agent Trajectories}
\label{sec:analysis:behavior}

% TODO TAB 7: trajectory-level metrics
% TODO 实验: parse SWE-agent log 提取 metrics


% -----------------------------------------------------------------------------
% 5.2 Failure Mode Analysis  (~ 0.3 页)
% -----------------------------------------------------------------------------
% FIG 8 [建议]: stacked bar chart, 失败类型分布对比
%
% 写作 takeaway: 提升集中在 "需要 function-level reasoning across files"
% 的任务, 对应 motivation.
% -----------------------------------------------------------------------------
\subsection{Where Does Mid-Training Help Most?}
\label{sec:analysis:failure}

% TODO FIG 8: failure mode breakdown
% TODO 分析: 标注 SWE-Bench task 类型


% -----------------------------------------------------------------------------
% 5.3 General Capability Forgetting Check  (~ 0.2 页)  [新增]
% -----------------------------------------------------------------------------
% Mid-training 是否 break 通用能力?
%   - MMLU / MATH / IFEval 在 14B 上 baseline vs +FIM 对比
%   - 如有小幅退化, 主动 report 并在 limitation 中 acknowledge
%
% TAB 8 [必须]: 通用 benchmark forgetting check
%
% CAUTION: 不测的话 reviewer 必然假设有 forgetting 问题.
%          主动 report 比被发现强.
% -----------------------------------------------------------------------------
\subsection{Preservation of General Capabilities}
\label{sec:analysis:forgetting}

% TODO TAB 8: forgetting check
% TODO 实验 (强烈建议): MMLU / MATH / IFEval 14B baseline vs +FIM


% =============================================================================
% 6. LIMITATIONS AND DISCUSSION  (~ 0.4 页)
% =============================================================================
% 必须 acknowledge:
%   1. Python-heavy benchmark; 跨语言 generalization 未充分验证.
%   2. CoT 监督依赖 frontier model API.
%   3. PDG 分析对动态语言不完美 (动态 dispatch, decorator, monkey patching).
%   4. Cross-base-model 只在 Qwen3-8B + SWE-Bench 上做, 未覆盖 Llama,
%      DeepSeek 等其他 family.
%   5. CoT-source ablation 仅在 14B 做 (compute 限制).
%   6. Function-aware FIM 假设代码 modular.
% =============================================================================
\section{Limitations and Discussion}
\label{sec:limitations}

% TODO: 6 条 limitations


% =============================================================================
% 7. CONCLUSION  (~ 0.25 页)
% =============================================================================
% 4-5 句:
%   - 重申 framing (FIM ↔ agent inductive bias)
%   - 重申 method
%   - 重申三轴 robustness 和 cross-domain transfer
%   - future direction
% =============================================================================
\section{Conclusion}
\label{sec:conclusion}

% TODO: 4-5 句 conclusion


% =============================================================================
% References & Appendix
% =============================================================================
\bibliographystyle{plainnat}
% \bibliography{refs}

\section*{References}
% TODO: 整理参考文献
% 必引文献 cluster:
%   FIM:           Bavarian 2022, Code-Llama, StarCoder, Qwen-Coder,
%                  DeepSeek-Coder
%   Mid-training:  MiniCPM, OLMo, 任何 explicit "mid-training" paper
%   Coding agent:  SWE-Bench, SWE-agent, OpenHands, Aider, r2e-gym, swe-smith,
%                  R2E, SWE-Gym
%   Distillation:  Orca, Self-Instruct, Distill-Whisper
%   Tool use eval: tau-bench, BFCL
%   General eval:  MMLU, MATH, IFEval (forgetting check)
%   Program analysis: PDG classical refs

\appendix

% =============================================================================
% APPENDIX
%
%   A. Detailed dataset statistics
%   B. Hyperparameters (含 LR sweep 详细数据)
%   C. Algorithm details
%      - PDG 实现细节 (Python AST 处理动态特性的 workaround)
%      - IT scoring reference model 选择
%      - CoT prompt template (Gemini-3 用什么 prompt)
%      - 14B-self CoT prompt template (与 Gemini-3 一致, 仅换模型)
%   D. Additional results
%      - 每个 benchmark 的 per-task breakdown
%      - Mid-train only (no post-train) sanity check 的结果
%      - Qwen3-8B 上更详细的 trajectory 分析 (如有)
%   E. Qualitative examples
%   F. Reproducibility checklist
% =============================================================================
\section{Dataset Details}
\label{app:data}

\section{Training Hyperparameters}
\label{app:hp}

\section{Algorithm Details}
\label{app:algo}

\section{Additional Results}
\label{app:results}

\section{Qualitative Examples}
\label{app:examples}

\input{checklist.tex}

\end{document}